
\subsection*{7. Appendice~: Th\'eor\`eme de l'indice de Hodge}

On retrouve et compl\`ete ici quelques r\'esultats ant\'erieurs par un raisonnement direct.

\textbf{Th\'eor\`eme 7.1.} \emph{Soit $X$ une surface irr\'eductible, propre sur un corps alg\'ebriquement clos. On suppose qu'il existe un faisceau inversible $H$ sur $X$ tel que $\langle c_1(H)^2\rangle>0$ (par exemple, $H$ ample). Alors un faisceau inversible $L$ sur $X$ tel que}
\[
   \langle c_1(L)\cdot c_1(H)\rangle=0
   \quad\hbox{\uline{et}}\quad
   \langle c_1(L)^2\rangle\geq 0,
\]
\emph{est num\'eriquement \'equivalent \`a $0$. (Par cons\'equent, la forme d'intersection sur $A^1_{\mathrm{num}}(X)_{\Q}$ ($=\bigl[\Pic(X)/\Pic^\tau(X)\bigr]\otimes\Q$, d'apr\`es 4.6) est non-d\'eg\'en\'er\'ee de type $(1,\rho-1)$).}

\textbf{7.1.1.} D'abord supposons : $X$ r\'eduit (donc, int\`egre) ; $H$ ample ; $\langle c_1(L)\cdot c_1(H)\rangle=0$ ; et $\langle c_1(L)^2\rangle>0$. Raisonnant par l'absurde, prenons $m>0$ tel que $H_1=L\otimes H^{\otimes m}$ soit ample, et prenons $n$ tel que, en posant $N=L^{\otimes n}\otimes H^{-1}$, on ait que $\langle c_1(N)\cdot c_1(H_1)\rangle=n\langle c_1(L)^2\rangle-m\langle c_1(H)^2\rangle$ est $>0$. Alors on a aussi que $\langle c_1(N)\cdot c_1(H)\rangle=-\langle c_1(H)^2\rangle$ est $<0$, et que $\langle c_1(N)^2\rangle=n^2\langle c_1(L)^2\rangle+\langle c_1(H)^2\rangle$ est $>0$, ce qui contredit le lemme suivant.

\textbf{Lemme 7.1.2.} \emph{Soit $X$ une surface int\`egre, projective sur un corps alg\'ebriquement clos. Pour un faisceau inversible $N$ sur $X$ tel que $\langle c_1(N)^2\rangle>0$, les conditions suivantes sont \'equivalentes}~:
\begin{enumerate}[label=(\roman*)]
\item \emph{Pour tout faisceau ample $H$, on a $\langle c_1(N)\cdot c_1(H)\rangle>0$.}
\item[(i')] \emph{Il existe un faisceau ample $H$ tel qu'on ait $\langle c_1(N)\cdot c_1(H)\rangle>0$.}
\item \emph{Il existe un entier $n>0$ tel qu'on ait $H^0(X,N^{\otimes n})\ne0$.}
\end{enumerate}

